\section{TODO: Unimplemented Features \& Roadmap}
\label{sec:todo}


This document tracks planned improvements and currently missing features for \texttt{xeus-haskell}.
\subsection{Core Kernel Features}


\begin{itemize}
\item [x] \textbf{\texttt{is\_complete} implementation}: Map parsability of mixed cells to Jupyter's completeness status.
\item [ ] \textbf{Incomplete Status}: Refine \texttt{is\_complete} to distinguish between \texttt{invalid} and \texttt{incomplete} code (e.g., unclosed delimiters).
\item [x] \textbf{\texttt{inspect\_request}}: Support for "Introspection" (Shift+Tab in Jupyter) to show documentation or type signatures for identifiers.
\item [ ] \textbf{\texttt{history\_request}}: Implementation of the Jupyter history protocol to allow searching and retrieving previous cell inputs.
\item [ ] \textbf{Advanced Completion}: Improve \texttt{completion\_request} to support qualified names (e.g., \texttt{Prelude.putStrLn}) and type-aware suggestions.
\item [ ] \textbf{Kernel Interrupt}: Add support for interrupting long-running Haskell executions (SIGINT handling).
\end{itemize}


\subsection{REPL Improvements}


\begin{itemize}
\item [ ] \textbf{Selective Re-compilation}: Instead of concatenating all previous definitions for every new definition, implement a dependency-tracking system to only re-compile what is necessary.
\item [ ] \textbf{Better Error Reporting}: Map MicroHs compile errors back to the specific line numbers in the Jupyter cell.
\item [ ] \textbf{WASM Optimization}: Reduce the size of the WebAssembly bundle and optimize the initial "warmup" time in the browser.
\item [x] \textbf{Language Extensions}: Work with the MicroHs upstream to support more modern Haskell extensions (e.g., \texttt{GADTs}, \texttt{TypeFamilies}).
\end{itemize}


\subsection{Rich Output Display System}


\begin{itemize}
\item [ ] \textbf{Standard Library Instances}: Add \texttt{Display} instances for more types in the MicroHs standard library.
\item [ ] \textbf{Data Visualization}: Create bridges for common Haskell charting libraries to output Jupyter-compatible JSON/Vega-Lite data.
\item [ ] \textbf{Widgets}: Initial investigation into supporting Jupyter Widgets (ipywidgets protocol).
\end{itemize}


\subsection{Infrastructure \& Testing}


\begin{itemize}
\item [ ] \textbf{Enhanced CI}: Add automated UI tests using \texttt{playwright} or \texttt{selenium} to verify the kernel works in a real JupyterLab/JupyterLite instance.
\item [ ] \textbf{Benchmarking}: Create a suite of performance benchmarks to track compilation and execution speed regressions.
\item [ ] \textbf{Conda-forge/Emscripten-forge}: Package the kernel for easier installation without needing to build from source.
\end{itemize}
